%% ma: L12-B08-0158
%% lop: 12
%% chuong: 2
%% bai: 8
%% dang: 12.8.51
%% loai: TN4
%% muc: VD
%% dap_an: "B"
%% nguon: "CĐ2.4.pdf giáo viên gửi 27/10/2026 (M07, Lê Minh Khôi – Nguyễn Anh Khôi, Chương II Chủ đề 4 Các bài toán về tâm tỉ cự), Câu 8"
%% kien_thuc: "Bài 8 (biểu thức tọa độ của phép toán vectơ, độ dài vectơ), Bài 6 (điểm thỏa hệ thức vectơ, quy tắc ba điểm), Bài 7 (hình chiếu lên mặt phẳng tọa độ)"
%% trang_thai: chinh-thuc
%% ghi_chu: "Giáo viên đã duyệt 27/10/2026"
\begin{ex}
Trong không gian $Oxyz$, cho ba điểm $A(-1;2;1)$, $B(2;-1;3)$, $C(3;5;-1)$. Điểm $M(a;b;c)$ thuộc mặt phẳng $(Oyz)$ thỏa mãn $\left|\overrightarrow{MA}+2\overrightarrow{MB}-\overrightarrow{CM}\right|$ đạt giá trị nhỏ nhất. Tính giá trị của biểu thức $E=2b+c$.
\choice{$E=-1$}{\True $E=4$}{$E=1$}{$E=-4$}
\loigiai{$-\overrightarrow{CM}=\overrightarrow{MC}$. Gọi $I$ thỏa $\overrightarrow{IA}+2\overrightarrow{IB}+\overrightarrow{IC}=\vec 0\Rightarrow I=\dfrac{A+2B+C}4=\left(\dfrac32;\dfrac54;\dfrac32\right)$.

Biểu thức bằng $4MI$, nhỏ nhất khi $M\left(0;\dfrac54;\dfrac32\right)$. Vậy $E=\dfrac52+\dfrac32=4$.}
\end{ex}
